\DocumentMetadata{
  pdfversion=2.0,
  pdfstandard=ua-2,
  lang=en-US,
  tagging=on,
  tagging-setup={math/setup=mathml-SE,tabsorder=structure}
}
\documentclass[11pt,letterpaper]{article}
\usepackage[margin=0.68in,includefoot]{geometry}
\usepackage{newcomputermodern}
\usepackage{amsmath,amscd,mathtools}
\usepackage{tikz,adjustbox}
\providecommand{\square}{\mdlgwhtsquare}
\usepackage[hidelinks]{hyperref}
\hypersetup{
  pdftitle={Algebra First-Year Exam, May 2013, Part 2},
  pdfauthor={University of Florida Department of Mathematics},
  pdfsubject={Graduate mathematics examination},
  pdfdisplaydoctitle=true,
  bookmarksopen=true
}
\setcounter{secnumdepth}{0}
\setlength{\parindent}{0pt}
\setlength{\parskip}{0.28em}
\setlength{\emergencystretch}{3em}
\setlength{\leftmargini}{2.15em}
\setlength{\leftmarginii}{2.65em}
\setlength{\leftmarginiii}{3em}
\pagestyle{empty}
\raggedbottom
\allowdisplaybreaks
\NewTaggingSocketPlug{block/list/label}{examproblem}{%
  \tagstructbegin{tag=Lbl}\tagstructend%
  \tagstructbegin{tag=\LBody}%
  \tagstructbegin{tag=\ProblemHeadingTag,title-o={\ProblemHeadingText}}%
  \tagmcbegin{tag=\ProblemHeadingTag,actualtext={\ProblemHeadingText}}%
  #2%
  \tagmcend%
  \tagstructend%
}
\newenvironment{examproblems}[2]{%
  \def\ProblemHeadingTag{#2}%
  \AssignTaggingSocketPlug{block/list/label}{examproblem}%
  \begin{enumerate}%
  \setcounter{enumi}{#1}%
  \renewcommand{\labelenumi}{\textbf{\theenumi.}}%
  \setlength{\topsep}{0.35em}%
  \setlength{\partopsep}{0pt}%
  \setlength{\itemsep}{0.48em}%
  \setlength{\parsep}{0.18em}%
}{\end{enumerate}}
\newenvironment{examsubparts}{%
  \AssignTaggingSocketPlug{block/list/label}{default}%
  \begin{enumerate}%
  \setlength{\topsep}{0.18em}%
  \setlength{\partopsep}{0pt}%
  \setlength{\itemsep}{0.16em}%
  \setlength{\parsep}{0.1em}%
}{\end{enumerate}}
\newenvironment{examinstructions}{%
  \par\begingroup\itshape
}{%
  \par\endgroup\vspace{0.18em}
}
\makeatletter
\renewcommand\subsection{\@startsection{subsection}{2}{0pt}%
  {-1.25ex plus -.25ex minus -.1ex}%
  {0.55ex plus .1ex}%
  {\normalfont\large\bfseries}}
\renewcommand{\@maketitle}{%
  \newpage\null\vskip -1em%
  {\footnotesize\bfseries
    \href{https://www.ufl.edu/}{University of Florida}, \href{https://math.ufl.edu/}{Department of Mathematics}\par}%
  \vskip -0.5em%
  \tagstructbegin{tag=H1,title={Algebra First-Year Exam, May 2013, Part 2}}%
  \tagpdfparaOff%
  \tagmcbegin{tag=H1}%
    {\LARGE\bfseries\href{https://gma.math.ufl.edu/exams/algebra-first-year/}{Algebra First-Year Exam}, May 2013, Part 2\par}%
  \tagmcend%
  \tagpdfparaOn%
  \tagstructend%
  \vskip 0.25em%
  \tagmcbegin{artifact}%
    \hrule height 1pt%
  \tagmcend%
  \vskip 1em%
}
\makeatother
\title{Algebra First-Year Exam, May 2013, Part 2}
\author{}
\date{}
\begin{document}
\maketitle
\thispagestyle{empty}
\begin{examinstructions}
Answer four problems. Write your answers clearly in complete English sentences. You may quote results (within reason) as long as you state them clearly.
\end{examinstructions}
\begin{examproblems}{0}{H2}
\def\ProblemHeadingText{Problem 1.}
\item
Let~\(R\) be an integral domain, let \(P \subset R\) be a prime ideal, and let \(F=\operatorname{Frac}(R)\) be the field of fractions of~\(R\). Define
\[
R_P=\left\{\frac{a}{b}:a,b\in R,\ b\notin P\right\}.
\]
\begin{examsubparts}
\item[\textnormal{(a)}]
Prove that \(R_P\) is a subring of~\(F\).
\item[\textnormal{(b)}]
Prove that \(R_P\) has a unique maximal ideal.
\end{examsubparts}
\def\ProblemHeadingText{Problem 2.}
\item
Let \(\omega=-\frac12+\frac12\sqrt{-3}\in\mathbb{C}\).
\begin{examsubparts}
\item[\textnormal{(a)}]
Prove that~\(\omega\) is a root of the polynomial \(f(X)=X^2+X+1\).
\item[\textnormal{(b)}]
Prove that the set \(\mathbb{Z}[\omega]=\{a+b\omega:a,b\in\mathbb{Z}\}\) is a subring of~\(\mathbb{C}\).
\item[\textnormal{(c)}]
Prove that \(\mathbb{Z}[\omega]\) is a Euclidean domain with respect to the norm \(N(\alpha)=\alpha\overline{\alpha}\), where \(\overline{\alpha}\) denotes the complex conjugate of \(\alpha\in\mathbb{Z}[\omega]\).
\end{examsubparts}
\def\ProblemHeadingText{Problem 3.}
\item
Let~\(R\) be an integral domain and let \(P\subset R\) be a prime ideal. Let
\[
f(X)=X^n+a_{n-1}X^{n-1}+\cdots+a_1X+a_0
\]
 be a monic polynomial of degree \(n\geq 1\) with coefficients in~\(R\).
\begin{examsubparts}
\item[\textnormal{(a)}]
Prove Eisenstein’s criterion: If \(a_i\in P\) for \(0\leq i\leq n-1\) and \(a_0\notin P^2\), then \(f(X)\) is irreducible in \(R[X]\).
\item[\textnormal{(b)}]
Give an example of \(f(X)\) as above such that \(a_i\in P\) for \(0\leq i\leq n-1\) but \(f(X)\) is not irreducible.
\end{examsubparts}
\def\ProblemHeadingText{Problem 4.}
\item
This problem has three parts.
\begin{examsubparts}
\item[\textnormal{(a)}]
Let~\(R\) be an integral domain and let~\(M\) be an~\(R\)-module. Define the rank of~\(M\).
\item[\textnormal{(b)}]
Prove that~\(\mathbb{Q}\) has rank 1 as a~\(\mathbb{Z}\)-module.
\item[\textnormal{(c)}]
Prove that~\(\mathbb{Q}\) is not a free~\(\mathbb{Z}\)-module.
\end{examsubparts}
\def\ProblemHeadingText{Problem 5.}
\item
Find a representative for each similarity class of noninvertible \(3\times 3\) matrices with entries in the field \(\mathbb{F}_2\) with 2 elements.
\def\ProblemHeadingText{Problem 6.}
\item
Let \(E/F\) be a field extension and let~\(R\) be a subring of~\(E\) which contains~\(F\).
\begin{examsubparts}
\item[\textnormal{(a)}]
Prove that if \(E/F\) is an algebraic extension then~\(R\) is a field.
\item[\textnormal{(b)}]
Give an example where \(E/F\) is not an algebraic extension and~\(R\) is not a field.
\end{examsubparts}
\end{examproblems}
\end{document}
